\chapter{Wie funktioniert ein Kampf?}
\label{chap:combat}

Ein Kampf ist eine besondere Szene, in der die Zeit in Runden und Züge aufgeteilt wird.

\section{Die Kampfrunde}

In jeder Runde handeln zuerst alle Helden, danach alle Gegner. Die Helden bestimmen ihre Reihenfolge selbst. Sobald alle Gegner vollständig aktiviert wurden, endet die Runde und die nächste beginnt.

\begin{kurzregel}[title=Reihenfolge einer Kampfrunde]
\textbf{Helden handeln} $\rightarrow$ \textbf{Gegner handeln}
\end{kurzregel}

\subsection{Der Zug eines Helden}

Ein Held besitzt in seinem Zug normalerweise:

\begin{itemize}
  \item eine \textbf{Standardaktion (\standardAction)},
  \item eine \textbf{Bewegungsaktion (\moveAction)},
  \item beliebig viele \textbf{freie Aktionen (\freeAction)}.
\end{itemize}

Mit einer \standardAction{} kann ein Held beispielsweise angreifen, einen Zauber wirken oder eine besondere Fähigkeit einsetzen.

Mit einer \moveAction{} kann er sich bewegen oder seine Ausrüstung neu anordnen. Seine \standardAction{} darf er stattdessen als zweite \moveAction{} verwenden.

\freeAction{} sind kurze Handlungen wie etwas rufen, auf etwas zeigen oder einen Gegenstand fallen lassen. Sie dürfen keine umfangreiche Handlung ersetzen.

Wird eine Aktion beschleunigt, wird eine \standardAction{} zur \moveAction{} und eine \moveAction{} zur \freeAction. Eine \freeAction{} kann nicht weiter beschleunigt werden.

\begin{tip}[title=Tipp: Aktionen zuerst planen]
Kläre zu Beginn deines Zuges kurz, was du mit deiner Standardaktion und deiner Bewegungsaktion erreichen möchtest. Das beschleunigt den Zug und verhindert unnötige Rückfragen während der Ausführung.
\end{tip}

\section{Bewegung, Nähe und Entfernung}

Figuren können zueinander \textbf{benachbart} oder \textbf{entfernt} sein. Eine Figur kann gleichzeitig zu mehreren Figuren benachbart sein.

Wie viele Bewegungen nötig oder mit einer \moveAction{} möglich sind, ergibt sich aus der Battlemap oder aus der Beschreibung der Spielleitung. Hindernisse und die konkrete Situation können eine Bewegung erschweren oder mehrere Bewegungen erforderlich machen.

Ein Held darf sich von einem benachbarten Gegner entfernen, ohne dass dies automatisch eine weitere Folge auslöst. Gegnerfähigkeiten können davon abweichen.

\begin{todo}
\textbf{Regeldesign: Bewegung im Kampf}

Die abstrakte Bewegung über \emph{benachbart} und \emph{entfernt} muss später mit den Battlemap-Regeln abgestimmt werden. Zu klären ist insbesondere, wie viele Felder oder Zonen eine Bewegungsaktion umfasst und wie Hindernisse behandelt werden.
\end{todo}

\section{Ausrüstung im Kampf wechseln}

Mit einer \moveAction{} darf ein Held seine ausgerüsteten und verstauten Gegenstände beliebig neu anordnen, solange danach alle Ausrüstungsplätze korrekt belegt sind.

Der Gegenstand am Körper-Ausrüstungsplatz ist davon ausgenommen. Eine Körperrüstung anzulegen, abzulegen oder zu wechseln, benötigt eine eigene \standardAction. Die allgemeinen Regeln für Ausrüstungsplätze und das Wechseln von Ausrüstung stehen in \chapterref{chap:equipment}.

\section{Einen Gegner angreifen}

Ein Angriff benötigt eine \standardAction. \textbf{Angriffsoptionen} beschreiben, wie eine Figur einen Angriff ausführt. Sie befinden sich meist auf \textbf{Waffen}, können aber auch auf \textbf{Zaubern} stehen. Vereinzelt verleihen auch andere \textbf{Ausrüstungsgegenstände} oder \textbf{Talente} eigene Angriffsoptionen. Die zugehörigen Regeln stehen in \chapterref{chap:equipment}, \chapterref{chap:magic} und \chapterref{chap:talents}.

Jede Angriffsoption legt fest,

\begin{itemize}
  \item welche Angriffsfertigkeit verwendet wird,
  \item welche Kaufoptionen zur Verfügung stehen und
  \item welche zusätzlichen Folgen nach dem Angriff eintreten.
\end{itemize}

Eine Angriffsoption auf einer Ausrüstungskarte kann nur gewählt werden, wenn sie auf der aktuell sichtbaren Kartenseite abgebildet ist.

\begin{kurzregel}[title=Einen Angriff durchführen]
\begin{enumerate}
  \item Angriffsoption wählen.
  \item Zulässiges Ziel bestimmen.
  \item Angriffsprobe ablegen.
  \item Bei mindestens einem \star{} Kaufoptionen erwerben.
  \item Verbleibende \star{} zum Verbessern erhaltener Schadenswürfel ausgeben.
  \item Schadenspools durch Effekte verändern und gegebenenfalls vor dem Wurf auf mehrere Ziele verteilen.
  \item Jeden Schadenspool getrennt werfen und den Schaden addieren.
  \item Rüstungsbruch und \armor{} für jedes Ziel anwenden.
  \item Verbleibenden Schaden von den jeweiligen Lebenspunkten abziehen.
  \item Gekaufte Kampfeffekte und zusätzliche Folgen zu ihrem jeweils angegebenen Zeitpunkt abhandeln.
\end{enumerate}
\end{kurzregel}

Für die Angriffsfertigkeiten gelten folgende Reichweiten:

\begin{itemize}
  \item \textbf{Schlagen} kann benachbarte Gegner als Ziel haben.
  \item \textbf{Schießen} kann entfernte Gegner als Ziel haben. Ein Held darf auch schießen, während ein Gegner zu ihm benachbart ist; das Ziel des Angriffs muss dennoch entfernt sein.
  \item \textbf{Zaubern} kann benachbarte oder entfernte Gegner als Ziel haben, sofern der Zauber nichts anderes bestimmt.
\end{itemize}

\subsection{Treffen}

Der Spieler legt mit der angegebenen Angriffsfertigkeit eine Angriffsprobe nach den Regeln aus \chapterref{chap:checks} ab. Zeigt sie keinen \star, verfehlt der Angriff. Jeder erzielte \star{} kann anschließend ausgegeben werden, um eine oder mehrere Kaufoptionen der gewählten Angriffsoption zu erwerben. Die vollständigen Regeln für Kaufoptionen stehen in \chapterref{chap:equipment}.

Helden erhalten Vorteil auf Angriffsproben gegen Gegner mit dem Statuseffekt \emph{Zu Boden}. Ein Held mit dem Statuseffekt \emph{Zu Boden} hat Nachteil auf seine eigenen Angriffsproben.

Gegner weichen Angriffen von Helden grundsätzlich nicht aus.

\subsection{Kaufoptionen und Schaden}

Es gibt drei Arten von Schadenswürfeln:

\begin{itemize}
  \item den \textbf{orangen Schadenswürfel (\orangeDmgDie)} für leichten Schaden,
  \item den \textbf{roten Schadenswürfel (\redDmgDie)} für mittleren Schaden,
  \item den \textbf{schwarzen Schadenswürfel (\blackDmgDie)} für schweren Schaden.
\end{itemize}

Eine Kaufoption nennt ihre Kosten in \star{} und den dafür erhaltenen Effekt. Jede Kaufoption der gewählten Angriffsoption darf pro Angriff höchstens einmal gekauft werden. Der Angreifer darf mehrere unterschiedliche Kaufoptionen in beliebiger Kombination kaufen, solange er ihre Kosten bezahlt und alle Voraussetzungen erfüllt.

Nachdem alle gewünschten Kaufoptionen gekauft wurden, dürfen verbleibende \star{} bereits erhaltene Schadenswürfel verbessern. Jeder ausgegebene \star{} verbessert einen Würfel um eine Stufe: \orangeDmgDie{} wird zu \redDmgDie{}, \redDmgDie{} wird zu \blackDmgDie{}. Derselbe Würfel darf mehrfach verbessert werden; \blackDmgDie{} können nicht weiter verbessert werden. Danach verfallen noch übrige \star{}.

Alle gekauften und verbesserten Schadenswürfel bilden einen gemeinsamen Schadenspool. Sie werden gemeinsam geworfen und ihre Ergebnisse anschließend addiert. Kaufoptionen können statt Schadenswürfeln oder zusätzlich zu ihnen Kampfeffekte verleihen.

\begin{kurzregel}[title=Schaden bestimmen]
\begin{enumerate}
  \item Kaufoptionen jeweils höchstens einmal kaufen.
  \item Alle erhaltenen Schadenswürfel zu einem Schadenspool zusammenfügen.
  \item Verbleibende \star{} ausgeben, um Schadenswürfel stufenweise zu verbessern.
  \item Den Schadenspool durch weitere Effekte verändern.
  \item Falls ein Effekt wie Spalten dies erlaubt, die Würfel vor dem Wurf auf Ziele verteilen.
  \item Jeden Schadenspool getrennt werfen und den gewürfelten Schaden addieren.
  \item Rüstungsbruch und Schutz durch \armor{} für jedes Ziel anwenden.
  \item Verbleibenden Schaden als Verlust von Lebenspunkten anwenden.
  \item Weitere Kampfeffekte zu ihrem angegebenen Zeitpunkt abhandeln.
\end{enumerate}
\end{kurzregel}

\begin{beispiel}[title=Beispiel: Schadenswürfel kaufen und verbessern]
Ein Schwert besitzt die Kaufoptionen 2\star{} $\rightarrow$ \orangeDmgDie{}\orangeDmgDie{} und 3\star{} $\rightarrow$ \orangeDmgDie{}\orangeDmgDie{}\orangeDmgDie{}. Mit 6\star{} kauft ein Held beide Optionen je einmal. Für den verbleibenden \star{} verbessert er einen der fünf orangen Würfel zu einem roten Würfel. Sein Schadenspool besteht aus 1\redDmgDie{} und 4\orangeDmgDie{}.
\end{beispiel}

\begin{beispiel}[title=Beispiel: Schaden und Kampfeffekt kombinieren]
Eine Axt besitzt unter anderem die Kaufoptionen 2\star{} $\rightarrow$ \redDmgDie{} und 2\star{} $\rightarrow$ \effectArmorBreak{} 2. Mit 4\star{} kann ein Held beide Optionen kaufen. Der Angriff verursacht Schaden mit 1\redDmgDie{} und erhält zusätzlich 2 Rüstungsbruch. Beide Kaufoptionen sind damit für diesen Angriff verbraucht.
\end{beispiel}

Das \textbf{Schild-Symbol (\armor)} beschreibt Schutz durch Rüstung, Schilde, Talente oder andere Effekte. Jedes \armor{} verhindert 1 Schaden. \textbf{Rüstungsbruch (\armorPierce)} ignoriert die angegebene Anzahl an Schild-Symbolen des Ziels, unabhängig von ihrer Quelle.

\begin{kurzregel}[title=Schild-Symbole und Rüstungsbruch]
Jedes \armor{} verhindert \textbf{1 Schaden}. Jedes \armorPierce{} ignoriert \textbf{1 \armor{}} des Ziels.
\end{kurzregel}

Der verbleibende Schaden lässt die Figur entsprechend viele \textbf{Lebenspunkte verlieren (\lifeLoss)}.

\section{Wenn Gegner angreifen}

Gegner legen keine Angriffsprobe ab. Ihr Profil gibt unmittelbar die verwendeten Schadenswürfel an.

Bevor diese Würfel geworfen werden, darf der angegriffene Held gegebenenfalls \textbf{ausweichen (\evade)}. Welche Würfel er für seine Ausweichprobe verwendet, ergibt sich aus seiner Ausrüstung oder seinen Talenten.

\begin{kurzregel}[title=Ausweichen]
Ein Held kann nur ausweichen, wenn seine \textbf{Ausrüstung oder ein Talent} dies erlaubt.

Jeder \star{} der Ausweichprobe entfernt einen beliebigen Schadenswürfel des Gegners. Überschüssige \star{} verfallen. Entfernte Schadenswürfel werden nicht geworfen.
\end{kurzregel}

Ein Held mit dem Statuseffekt \emph{Zu Boden} hat Nachteil auf seine Ausweichprobe. Greift ein Gegner mit dem Statuseffekt \emph{Zu Boden} an, hat der angegriffene Held Vorteil auf seine Ausweichprobe. Kann der Held aufgrund seiner Ausrüstung oder Talente nicht ausweichen, entfällt die Probe vollständig.

Die verbleibenden Schadenswürfel werden geworfen. Anschließend verhindert jedes \armor{} des Helden 1 Schaden. Für den übrigen Schaden verliert der Held entsprechend viele Lebenspunkte \lifeLoss.

\begin{todo}
\textbf{Balancing: Wahl der Schadenswürfel beim Ausweichen}

Aktuell darf der Held bei einem gemischten Schadenspool frei wählen, welche Schadenswürfel durch seine Ausweicherfolge entfernt werden. Prüfen, ob diese freie Wahl dauerhaft sinnvoll und ausgewogen ist.
\end{todo}

\section{Kampfunfähig oder besiegt}

Lebenspunkte können niemals unter 0 sinken.

Sinkt ein Gegner auf 0 Lebenspunkte, ist er besiegt.

Sinkt ein Held auf 0 Lebenspunkte, ist er kampfunfähig und kann nicht handeln. Wird er geheilt und erhält mindestens 1 \textbf{Lebenspunkt (\lifeGain)}, ist er wieder kampfbereit, erhält aber den Statuseffekt \emph{Zu Boden}. Er muss daher zunächst entsprechend den Regeln in \appendixref{app:status-effects} aufstehen.

Wenn alle Helden kampfunfähig sind, sind die Helden besiegt. Das bedeutet nicht automatisch, dass sie sterben.

\section{Wann endet ein Kampf?}

Ein Kampf endet, wenn die gefährliche Auseinandersetzung vorbei ist, beispielsweise weil alle Gegner besiegt sind, eine Seite flieht oder aufgibt oder das Ziel des Kampfes erreicht wurde.

Wenn der Kampf endet, wird wieder normal weitererzählt.

\section{Beispiele}

\begin{beispiel}[title=Beispiel: Ein Held greift an]
\begin{enumerate}
  \item \textbf{Angriffsoption wählen (Zahnschwert):} Kiera greift einen benachbarten Gegner mit ihrem Zahnschwert an.
  \item \textbf{Angriffsprobe:} Sie legt die angegebene Schlagen-Probe ab und erzielt 6\star.
  \item \textbf{Kaufoptionen erwerben:} Kiera kauft für 2\star{} zwei \redDmgDie{} und für 3\star{} drei \orangeDmgDie{}.
  \item \textbf{Schadenswürfel verbessern:} Mit dem verbleibenden \star{} verbessert sie einen \orangeDmgDie{} zu einem \redDmgDie{}.
  \item \textbf{Schaden würfeln:} Ihr Schadenspool aus 3\redDmgDie{} und 2\orangeDmgDie{} verursacht zusammen 8 Schaden.
  \item \textbf{Schutz anwenden:} Der Gegner besitzt 2\armor{} und verhindert damit 2 Schaden.
  \item \textbf{Lebenspunkte verlieren:} Der Gegner verliert die verbleibenden 6 Lebenspunkte.
\end{enumerate}
\end{beispiel}

\begin{beispiel}[title=Beispiel: Ein Gegner greift an]
\begin{enumerate}
  \item \textbf{Schadenspool bestimmen:} Ein Gegner greift Kiera mit 1\blackDmgDie{} und 2\orangeDmgDie{} an.
  \item \textbf{Ausweichen:} Kieras Ausrüstung erlaubt eine Ausweichprobe. Sie erzielt 1\star{} und entfernt den \blackDmgDie.
  \item \textbf{Schaden würfeln:} Die beiden verbleibenden \orangeDmgDie{} verursachen zusammen 3 Schaden.
  \item \textbf{Schutz anwenden:} Kiera besitzt 1\armor{} und verhindert damit 1 Schaden.
  \item \textbf{Lebenspunkte verlieren:} Kiera verliert die verbleibenden 2 Lebenspunkte.
\end{enumerate}
\end{beispiel}
